\section{什么叫“资金先行、价格滞后”}

连续流入个股表提供了比一级行业更细的观察窗口。我们重点保留连续5日以上、累计涨幅不超过6\%的标的，再用一年价格位置、基本面和估值做二次筛选。这样可以剔除已经大涨的灵康药业、百合花，也可以识别电子行业内部的半导体设备资金孤岛。

\begin{exhibitbox}[Exhibit 5：连续流入且价格尚未明显反应的代表标的]
\begin{tabularx}{\textwidth}{L{2.3cm}L{2.3cm}R{1.6cm}R{2.0cm}R{1.8cm}X}
\toprule
\textbf{代码} & \textbf{名称} & \textbf{天数} & \textbf{净流入/亿元} & \textbf{同期涨跌} & \textbf{二次判断} \\
\midrule
601825 & 沪农商行 & 10 & 1.42 & -3.04\% & 低位、高股息、低杠杆；通过 \\
601077 & 渝农商行 & 8 & 1.46 & +3.11\% & 基本面改善最强；通过 \\
000425 & 徐工机械 & 5 & 2.60 & +0.12\% & 行业流出中的资金孤岛；通过 \\
601939 & 建设银行 & 7 & 3.91 & +4.81\% & 位置不低、弹性有限；观察 \\
601088 & 中国神华 & 5 & 1.87 & +3.29\% & 防御承接，周期与股息平衡；观察 \\
688125 & 安达智能 & 9 & 1.71 & +0.94\% & 价格回撤大但历史涨幅和估值仍高；不入核心 \\
002937 & 兴瑞科技 & 7 & 2.02 & -3.13\% & 资金吸纳但缺完整估值锚；观察 \\
300790 & 宇瞳光学 & 6 & 2.06 & -9.33\% & 下跌中承接，尚未确认趋势反转 \\
\bottomrule
\end{tabularx}
\end{exhibitbox}

\section{半导体设备：资金还在，低位已经不在}

中微公司连续6日净流入32.70亿元、同期上涨6.26\%；富创精密连续6日流入9.49亿元、同期涨4.81\%。如果只看这一周，它们仍像“资金先行”。但拉长到一年，7月10日收盘后，中微公司位于一年价格区间约90\%，富创精密约82\%；过去60日涨幅分别约105\%和184\%。因此，旧报告中“半导体设备全年横盘低位”的判断已经失效。

\begin{riskbox}[不能把回撤一天重新定义为低位]
中微公司和富创精密7月10日分别下跌7.86\%和9.30\%，但单日回撤并没有消除此前数月的大幅重估。当前更适合定义为高景气、高波动、等待估值消化，而不是“尚未启动”。
\end{riskbox}

\section{军工与传媒：已经越过潜伏线}

军工全周下跌仍获资金净流入，说明7月10日前确有提前承接；但商业航天事件推动板块当日大幅上涨，航天ETF还出现连续申购，阶段已经从左侧潜伏转为右侧确认。传媒周涨2.01\%、周流入27.33亿元、当日上涨2.60\%，同样属于早期启动，不适合再用“悄悄”描述。

\section{医药和机械：行业与个股必须分开}

医药7月10日净流入41.55亿元，但全周净流出53.45亿元；这更像政策催化后的回补。机械设备当日净流入16.57亿元，但全周净流出162.71亿元、周跌5.71\%，行业尚未形成持续承接。徐工机械之所以入选，不是因为机械设备行业获资金认可，而是因为它在行业失血时仍连续5日流入且价格几乎未涨，是典型的\textbf{资金孤岛}。
